\section{Related Work}
\label{sec:related_work}

Our methodological audit builds upon and critiques several intersecting fields of deep learning: static weight-space merging, adaptive test-time model merging, few-shot parameter tuning, and foundational reality-check audits.

\subsection{Weight-Space Model Merging}
Model merging operates on the principle that neural networks fine-tuned from a shared pre-trained initialization lie within a single, highly compatible loss basin \cite{izmailov2018averaging,cha2021swad}.
This linear mode connectivity allows weight-space interpolation to combine capabilities from disjoint experts without retraining.
Task Arithmetic \cite{ilharco2022editing} utilizes task vectors, formed by subtracting the pre-trained weights from fine-tuned weights, which are scaled and added back to the pre-trained model.
Model Soups \cite{wortsman2022model} averages weights from different fine-tuning runs to improve generalizability.
However, simple linear averaging often suffers from representational interference, where conflicting parameters cancel out task-specific features.
To alleviate this, Ties-Merge \cite{yadav2023ties} prunes redundant parameters and resolves sign conflicts before averaging, whilst Matena \& Raffel \cite{matena2021merging} employ diagonal Fisher information matrices to perform weight-space amalgamation.
Additionally, Jin et al. \cite{jin2022dataless} and Gu et al. \cite{gu2024advancing} explore data-free and parameter-efficient merging, respectively.
Despite these advancements, these baseline methods rely heavily on uniform, hand-tuned coefficients across all layers, which are suboptimal when layer-wise sensitivities diverge.

\subsection{Adaptive Model Merging \& Test-Time Adaptation}
To move beyond static uniform coefficients, adaptive model-merging methods optimize coefficients on a layer-wise basis.
Under the "no-data" constraint, these methods rely on online Test-Time Adaptation (TTA) over target evaluation streams.
Inspired by fully test-time adaptation frameworks like Tent \cite{wang2020tent,liang2023source}, AdaMerging \cite{adamerging} optimizes layer-wise, task-specific coefficients online by minimizing the Shannon prediction entropy of incoming batches via gradient descent.
RegCalMerge \cite{regcalmerge} extends this by introducing class-capacity normalization and elastic spatial regularization to prevent sacrificial task bias under class-imbalanced streams.
PolyMerge \cite{polymerge} restricts layer-wise coefficients to a low-degree polynomial function of depth to shrink the optimization space.
To mitigate the transductive stream noise inherent in online streams, traditional TTA pipelines often employ temporal smoothing, such as exponential moving averages (EMA) of model parameters \cite{wang2020tent}.
While online TTA is theoretically appealing, it relies on several ideal assumptions—such as balanced streams and noise-free updates—and requires additional test-time computation and backpropagation during deployment.
Crucially, as we show in our physical validation, even with temporal parameter smoothing, online unconstrained methods remain highly vulnerable to the fundamental misalignment of the unsupervised entropy objective.

\subsection{Offline Few-Shot Validation Tuning}
Recent work has questioned the "no-data" premise of test-time adaptation, arguing that it represents a "strawman" scenario \cite{liao2021are}.
In most physical engineering pipelines, practitioners have access to a very small, labeled validation set (e.g., $M=10$ labeled samples per task) representing the target domain.
This small-scale data can be leveraged offline via Offline Few-Shot Validation Tuning (OFS-Tune).
OFS-Tune directly optimizes the merging coefficient trajectories offline using gradient-free local search or backpropagation to minimize validation cross-entropy.
While global derivative-free optimization techniques like Bayesian Optimization (BO) represent highly standard solutions for low-dimensional parameter search, we show that simple local search algorithms like the Nelder-Mead simplex method are exceptionally well-suited for this regime, converging rapidly without Gaussian Process modeling overhead.
Because validation tuning is completed entirely offline, it eliminates the need for test-time backpropagation, reducing deployment latency and compute requirements to zero.

\subsection{Methodological Auditing in Machine Learning}
Our work is philosophically aligned with "reality check" papers that systematically audit popular deep learning research paradigms.
For instance, Musgrave et al. \cite{musgrave2020metric} conducted a rigorous audit of deep metric learning, revealing that standard SOTA claims were often artifacts of unfair hyperparameter tuning and weak baselines.
In domain generalization, Gulrajani \& Lopez-Paz \cite{gulrajani2020in} showed that most proposed algorithms failed to outperform a simple empirical risk minimization (ERM) baseline when evaluated under a fair, standardized protocol.
Similarly, Liao et al. \cite{liao2021are} audited test-time adaptation, highlighting how un-regularized online adaptation overfits to stream statistics and collapses under out-of-distribution shifts.
By applying this methodology-driven auditing lens to the adaptive model-merging literature, we expose how the standard 4-task visual suite serves as a major confounding variable, masking severe transductive overfitting and representation collapse in unconstrained online TTA.
